\section{Invariants of monoidal categories}
\label{Sec : invariants of monoidal categories}
\subsection{Settings in this paper}\label{ssec : settings in this paper}
From now on, we take one of the following choices
for a monoidal category $\mathscr{C}$, an integer $N$,
modules $S_k$ and $T_k^{t_0}$, a matrix $A$, and a ring $K(\mathscr{C})_{\Z}$:
\begin{enumerate}
    \item[(i)] Let $\mathfrak{g}$ be a finite-dimensional simple Lie algebra
    and $\mathscr{C}_{\mathfrak{g}}$ be the category
    of finite-dimensional representations of $U_q(L\mathfrak{g})$. 
    Let $A = (a_{\iota \jmath})_{\iota, \jmath \in I_{\mathrm{fin}}}$ 
    be the Cartan matrix of $\mathfrak{g}_{\mathrm{fin}}$. 
    Let $[a, b]$ be an interval in $\Z$ of length $N$
    and $\mathfrak{s} = ((\iota_k, p_k))_{k \in \Z}$
    be an admissible sequence.
    Let $\mathscr{C}$ be the monoidal category
    $\mathscr{C}_{\mathfrak{g}}^{[a, b], \mathfrak{s}}$.
    Replacing $\mathfrak{s}$ by the admissible sequence
    $((\iota_{k + a - 1}, p_{k + a - 1}))_{k \in \Z}$,
    we may assume without loss of generality that $[a, b] = [1, N]$.
    Let $S_k = S_k^{\mathfrak{s}}$ be the cuspidal module
    in $\mathscr{C}_{\mathfrak{g}}^{[1, N], \mathfrak{s}}$
    for $1 \leq k \leq N$.
    Let $T_k^{t_0}$ be the simple module whose PBW parameter is
    $\sum_{\substack{1\leq s \leq k \\ \iota_s = \iota_k}} e_s$.
    Let
    \[
    K(\mathscr{C})_{\Z} \coloneq K(\mathscr{C})
    \simeq \Z\bigl[[S_k] \mid 1 \leq k \leq N\bigr]
    \]
    be the Grothendieck ring
    of $\mathscr{C}_{\mathfrak{g}}^{[1, N], \mathfrak{s}}$.
    \item[(ii)] Let $\mathfrak{g}$ be a symmetric Kac--Moody Lie algebra
    and $R$ be a symmetric quiver Hecke algebra
    corresponding to $\mathfrak{g}$
    over a base field of characteristic $0$. 
    Let $A = (a_{\iota, \jmath})$ be the Cartan matrix of $\mathfrak{g}$. 
    Let $w \in W$ and
    $\widetilde{w} = s_{\iota_1}\cdots s_{\iota_N}$
    be a reduced word.
    Let $\mathscr{C}$ be the monoidal category $\mathscr{C}_w$,
    and let $S_k$ be the cuspidal module for $1 \leq k \leq N$.
    Let $T_k^{t_0} \coloneq M[k^{\min}, k]$.
    Let
    \[
    K(\mathscr{C})_{\Z} \coloneq K(\mathscr{C}_w)|_{q = 1}
    \simeq \Z\bigl[[S_k] \mid 1 \leq k \leq N\bigr]
    \]
    be the Grothendieck ring of $\mathscr{C}_w$
    specialized at $q = 1$.
\end{enumerate}

In what follows, for $L, M \in \mathscr{C}$,
even if $\mathscr{C} = \mathscr{C}_w$,
we write the tensor product or the convolution product
of $L$ and $M$ as $L \otimes M$.
Also, we ignore the grading shifts when $\mathscr{C} = \mathscr{C}_w$.

\subsection{Invariants}
\begin{Def}
\label{Def : commutativity reality and primeness}
\begin{enumerate}
    \item[(i)] Let $L, M \in \Irr\mathscr{C}$.
    We say that $L$ and $M$ \emph{commute}
    if $L \otimes M \simeq M \otimes L$.
    We say that $L$ and $M$ \emph{strongly commute}
    if $L \otimes M \in \Irr\mathscr{C}$.
    \item[(ii)] Let $L \in \Irr\mathscr{C}$.
    We call $L$ \emph{real} if $L \otimes L \in \Irr\mathscr{C}$.
    We call $L$ \emph{imaginary}
    if $L \otimes L \notin \Irr\mathscr{C}$.
    We call $L$ \emph{prime} if $L$ is non-trivial and,
    if $L \simeq L_1 \otimes L_2$, then one of them is trivial.
\end{enumerate}
\end{Def}

Since $K(\mathscr{C})_{\Z}$ is commutative,
strong commutativity implies commutativity.

Let $\Omega$ be the smallest set of non-zero objects in $\mathscr{C}$
which contains $\Irr\mathscr{C}$
and is closed under taking non-zero subquotients and tensor products.
In \cite[$\S$3]{KKOP20} and \cite[$\S$3]{KKKO18},
the invariant $\Lambda \colon \Omega \times \Omega \to \Z$
is defined using $R$-matrices.
Also, the symmetric invariant
$\mathfrak{d} \colon \Omega \times \Omega \to \Z$ is defined by
\[
\mathfrak{d}(M, N)
\coloneq \frac{1}{2}\bigl(\Lambda(M, N) + \Lambda(N, M)\bigr).
\]
Actually, it is known that $\mathfrak{d}(M, N) \in \Z_{\geq 0}$
for $M, N \in \Omega$
\cite[$\S$3]{KKOP20}, \cite[$\S$3]{KKKO18}.
When $\mathscr{C} = \mathscr{C}_{\mathfrak{g}}^{[1, N], \mathfrak{s}}$,
the invariant $\Lambda^{\infty} \colon \Omega \times \Omega \to \Z$
is defined in \cite[Definition 3.6]{KKOP20}.
When $\mathscr{C} = \mathscr{C}_w$,
let $\Lambda^{\infty} \colon \Omega \times \Omega \to \Z$
be the invariant such that
\[
\Lambda^{\infty}(M, N)
= -(\operatorname{wt} M, \operatorname{wt} N).
\]
Let
\[
\widetilde{\Lambda}(M, N)
\coloneq \frac{1}{2}
\bigl(\Lambda(M, N) - \Lambda^{\infty}(M, N)\bigr).
\]
It is known that $\widetilde{\Lambda}(M, N) \in \Z$
\cite[Definition 3.6]{KKOP20}, \cite[Lemma 3.1.2]{KKKO18}.

\begin{Rem}
\label{Rem : convention for tilde Lambda}
This $\widetilde{\Lambda}$ is compatible with the definition
of $\widetilde{\Lambda}$ in \cite[Definition 3.1.1]{KKKO18}
for the quiver Hecke case.
However, in \cite[Definition 3.6]{KKOP20}
for the quantum affine case,
the original $\widetilde{\Lambda}$ is defined by
$\frac{1}{2}(\Lambda + \Lambda^{\infty})$.
Since the $\Lambda^{\infty}$-invariant for quantum affine algebras
corresponds to $-(\operatorname{wt}, \operatorname{wt})$
for quiver Hecke algebras via quantum Schur--Weyl duality
\cite[Theorem 4.12]{PBW24},
we define $\widetilde{\Lambda}$ for quantum affine algebras by
$\frac{1}{2}(\Lambda - \Lambda^{\infty})$.
\end{Rem}

\begin{Prop}[{\cite[Corollary 3.17]{KKOP20},
\cite[Lemma 3.2.3]{KKKO18}}]
\label{Prop : d vanishing and commutativity}
Let $L, M \in \Irr\mathscr{C}$ and suppose one of them is real.
Then the following are equivalent:
\begin{enumerate}
    \item[(i)] $\mathfrak{d}(L, M) = 0$,
    \item[(ii)] $L$ and $M$ commute,
    \item[(iii)] $L$ and $M$ strongly commute.
\end{enumerate}
\end{Prop}

\begin{Lem}[{\cite[Lemma 3.10, Proposition 3.11]{KKOP20},
\cite[Lemma 3.1.5, Proposition 3.2.10]{KKKO18}}]
\label{Lem : Lambda tensor products and subquotients}
Let $L, M, N \in \Irr\mathscr{C}$.
\begin{enumerate}
    \item[(i)] We have
    \[
    \begin{aligned}
    \Lambda(L, M \otimes N) &= \Lambda(L, M) + \Lambda(L, N), \\
    \Lambda(M \otimes N, L) &= \Lambda(M, L) + \Lambda(N, L).
    \end{aligned}
    \]
    \item[(ii)] Let $S$ be a simple subquotient of $M \otimes N$.
    Then
    \[
    \begin{aligned}
    \Lambda(L, S) &\leq \Lambda(L, M) + \Lambda(L, N), \\
    \Lambda(S, L) &\leq \Lambda(M, L) + \Lambda(N, L).
    \end{aligned}
    \]
\end{enumerate}
\end{Lem}

By Proposition~\ref{Prop : d vanishing and commutativity}
and Lemma~\ref{Lem : Lambda tensor products and subquotients},
if $L$ is real, then $L^{\otimes n}$ is real for every $n \in \Z_{\geq 0}$.

\begin{Prop}[{\cite[Proposition 4.2]{KKOP20},
\cite[Proposition 3.2.10]{KKKO18}}]
\label{Prop : commutativity with simple subquotients}
Let $L, M, N \in \Irr\mathscr{C}$.
When $L$ is real, the following conditions are equivalent:
\begin{enumerate}
    \item[(a)] $L$ strongly commutes with $M$ and $N$.
    \item[(b)] Any simple subquotient $S$ of $M \otimes N$
    commutes with $L$ and satisfies
    \[
    \Lambda(S, L) = \Lambda(M, L) + \Lambda(N, L).
    \]
    \item[(c)] Any simple subquotient $S$ of $M \otimes N$
    commutes with $L$ and satisfies
    \[
    \Lambda(L, S) = \Lambda(L, M) + \Lambda(L, N).
    \]
\end{enumerate}
\end{Prop}

\begin{Lem}[{\cite[Lemma 3.10, Proposition 3.11]{KKOP20}}]
\label{Lem : Lambda infinity tensor products and subquotients}
Let $L, M, N \in \Irr\mathscr{C}$.
Then
\[
\begin{aligned}
\Lambda^{\infty}(L, M \otimes N)
&= \Lambda^{\infty}(L, M) + \Lambda^{\infty}(L, N), \\
\Lambda^{\infty}(M \otimes N, L)
&= \Lambda^{\infty}(M, L) + \Lambda^{\infty}(N, L).
\end{aligned}
\]
Let $S$ be a simple subquotient of $M \otimes N$.
Then
\[
\begin{aligned}
\Lambda^{\infty}(L, S)
&= \Lambda^{\infty}(L, M) + \Lambda^{\infty}(L, N), \\
\Lambda^{\infty}(S, L)
&= \Lambda^{\infty}(M, L) + \Lambda^{\infty}(N, L).
\end{aligned}
\]
\end{Lem}

\begin{Prop}[{\cite[Theorems 3.2 and 3.12]{headsimplicity}}]
\label{Prop : simplicity of heads and socles}
Let $L, M \in \Irr\mathscr{C}$
and assume that at least one of them is real.
Then the head and the socle of $L \otimes M$ are simple.
\end{Prop}

We write the head of $L \otimes M$ as $L \nabla M$
and the socle as $L \Delta M$. 

\begin{Lem}[{\cite[Corollary~3.7, Corollary~3.14]{headsimplicity}}]
\label{Lem : hd soc injectivity}
  Let $L \in \Irr \mathscr{C}$ be a real module. 
  Then the maps 
  \[
  L \nabla - , L \Delta -: \Irr \mathscr{C} \to  \Irr \mathscr{C}
  \]
  are both injective. 
\end{Lem}

\begin{Lem}[{\cite[Theorem 4.11]{KKOP20},
\cite[Theorem 4.1.1]{KKKO18}}]
\label{Lem : Lambda inequalities for heads and socles}
Let $M, N \in \Irr\mathscr{C}$ and assume that $M$ is real.
\begin{enumerate}
    \item[(i)] $[M \otimes N] = [M \nabla N] + \sum_k[S_k]$
    with simple modules $S_k$ such that
    \[
    \Lambda(M, S_k) < \Lambda(M, M \nabla N) = \Lambda(M, N);
    \]
    \item[(ii)] $[M \otimes N] = [M \Delta N] + \sum_k[S_k]$
    with simple modules $S_k$ such that
    \[
    \Lambda(S_k, M) < \Lambda(M \Delta N, M) = \Lambda(N, M);
    \]
    \item[(iii)] $[N \otimes M] = [N \nabla M] + \sum_k[S_k]$
    with simple modules $S_k$ such that
    \[
    \Lambda(S_k, M) < \Lambda(N \nabla M, M) = \Lambda(N, M);
    \]
    \item[(iv)] $[N \otimes M] = [N \Delta M] + \sum_k[S_k]$
    with simple modules $S_k$ such that
    \[
    \Lambda(M, S_k) < \Lambda(M, N \Delta M) = \Lambda(M, N).
    \]
\end{enumerate}
In particular, $M \nabla N$ as well as $M \Delta N$
appears only once in the Jordan--H\"older series
of $M \otimes N$ in $\mathscr{C}$.
\end{Lem}

\begin{Lem}[{\cite[Corollary 3.18]{LocalizationsI},
\cite[Proposition 2.17]{PBW24}}]
\label{Lem : decrease of d}
Let $L, M \in \Irr\mathscr{C}$ and assume that $L$ is real.
\begin{enumerate}
    \item[(i)] If $\mathfrak{d}(L, M) > 0$,
    then for any simple subquotient $S$ of $L \otimes M$,
    \[
    \mathfrak{d}(L, S) < \mathfrak{d}(L, M).
    \]
    \item[(ii)] There exists $n \in \Z_{\geq 0}$ such that
    \[
    \mathfrak{d}(L, S) = 0
    \]
    for any simple subquotient $S$ of $L^{\otimes n}\otimes M$. 
\end{enumerate}
\end{Lem}

\begin{Lem}[{\cite[Proposition 4.7]{KKOP20},
\cite[Proposition 3.2.17]{KKKO18}, \cite[Lemma 2.28]{PBW24}}]
\label{Lem : d one exact sequence}
Let $L, M \in \Irr\mathscr{C}$ and assume that one of them is real.
Assume that $\mathfrak{d}(L, M) = 1$.
Then
\[
0 \longrightarrow L \Delta M
\longrightarrow L \otimes M
\longrightarrow L \nabla M \longrightarrow 0
\]
is an exact sequence. 
In addition, if $L, M$ are both real, then $L \Delta M$ and $L \nabla M$ are real modules.
\end{Lem}

\begin{Lem}[{\cite[Lemma 2.23]{KKOP24}}]
\label{Lem : reality of head commuting with a factor}
Let $L, M \in \Irr\mathscr{C}$ be real modules.
Assume that $\mathfrak{d}(L, L \nabla M) = 0$.
Then $L \nabla M$ is real.
\end{Lem}

\begin{Prop}[{\cite[Theorem 4.13]{KKOP20},
\cite[Theorem 4.1]{KK19}}]
\label{Prop : divisibility by a real simple module}
Let $X, L \in \Irr\mathscr{C}$.
Assume that $L$ is real.
If $[X] = [L]\phi$ for some $\phi \in K(\mathscr{C})_{\Z}$,
then $X \simeq L \otimes Y$ for some $Y \in \Irr\mathscr{C}$
such that $\mathfrak{d}(L, Y) = 0$.
\end{Prop}

For $k \in [1, N]$, let $\beta_k \coloneq s_{\iota_1}\cdots s_{\iota_{k - 1}}\alpha_{\iota_k}$. 

\begin{Lem}[{\cite[Proposition 5.4, Proposition 5.11]{Laurentfamily}}]
\label{Lem : Lambda and PBW}
We set
\[
\lambda_{st} \coloneq (-1)^{\delta(s > t)}\delta(s \neq t)
(\beta_s, \beta_t)
\qquad \text{for }s, t \in [1, N].
\]
Here, for a proposition $P$,
\[
\delta(P) =
\begin{cases}
1 & \text{if }P\text{ is true}, \\
0 & \text{if }P\text{ is false}.
\end{cases}
\]
Let $L, M \in \Irr \mathscr{C}$. 
Assume that at least one of them is real and they commute. 
Then
\[
\Lambda(L, M) = \mathrm{PBW}(L)^{T}(\lambda_{kl})_{1 \leq k, l\leq N}\mathrm{PBW}(M). 
\]
\end{Lem}

\begin{Lem}[{\cite[Lemma 3.2]{simplylaced}}]
\label{Lem : Lambda infinity of cuspidal modules}
Let $k, l \in [1, N]$.
Then
\[
\Lambda^{\infty}(S_k, S_l) = -(\beta_k, \beta_l).
\]
\end{Lem}
For $\mathbf{a}, \mathbf{b} \in \Z^N$, we write $\mathbf{a}<_{\mathrm{lex}}\mathbf{b}$ by 
if $a_1=b_1, \ldots, a_{k-1}=b_{k-1}$ and $a_k<b_k$ for some $1 \leq k\leq N$. 
This is a total order. 

\begin{Lem}[{\cite[Theorem 6.12]{PBW24}, \cite[Proposition 4.31]{Kim12}}]\label{Lem : tensor PBW}
\begin{enumerate}
  \item Let $L \in \Irr \mathscr{C}$ and $\mathbf{a} = \mathrm{PBW}(M)$. 
  Then 
  \[
  [L] \in  \bigl[\bigotimes_{k=1}^N S_k^{\otimes a_k}\bigr] 
  + \sum_{\mathbf{b}<_{\mathrm{lex}}\mathbf{a}}\Z \bigl[\bigotimes_{k=1}^N S_k^{\otimes b_k}\bigr]
  \]
  \item Let $L_1, L_2 \in \Irr \mathscr{C}$ and $\mathbf{a}_1, \mathbf{a}_2$ be their PBW parameters. 
  Then 
  \[
  [L_1\otimes L_2] \in [L]+ \sum_{\mathrm{PBW}(M)<_{\mathrm{lex}}\mathbf{a}_1+\mathbf{a}_2}  \Z [M]
  \]
  in $K(\mathscr{C})_{\Z}$ where $\mathrm{PBW}(L) = \mathbf{a}_1 + \mathbf{a}_2$. 
\end{enumerate}
\end{Lem}

\begin{Prop}
\label{Prop : tensor factors in category}
Let $M,N$ be nonzero objects in $\mathscr{C}_{\mathfrak{g}}$ for a quantum affine case and 
in $R\text{-gmod}$ for a quiver Hecke case. 
If $M\otimes N\in\mathscr{C}$, then $M,N\in\mathscr{C}$. 
In particular, if $L \in \Irr\mathscr{C}$ is non-trivial and does not have non-trivial tensor factorization 
in $\mathscr{C}$, then $L$ is a prime module in the entire category 
$\mathscr{C}_{\mathfrak{g}}$ or $R\text{-gmod}$. 
\end{Prop}

\begin{proof}
First, we consider the case
$\mathscr{C}=\mathscr{C}_{\mathfrak{g}}^{[1, N], \mathfrak{s}}$.
Since an object belongs to $\mathscr{C}$ if and only if all its
composition factors belong to $\mathscr{C}$, we may assume that
$M,N\in\Irr \mathscr{C}_{\mathfrak{g}}$.
Let $m_1,m_2$ be the highest $\ell$-weight monomials of $M,N$. 
By Proposition~\ref{Prop : PBW parameterization}, the assumption
$M\otimes N\in\mathscr{C}$ implies that $m_1m_2$ is a monomial in
$\{Y_{\pi(\iota_k),p_k}\mid k\in[1,N]\}$.
Thus $m_1$ and $m_2$ are also monomials in
$\{Y_{\pi(\iota_k),p_k}\mid k\in[1,N]\}$, and hence $M,N\in\mathscr{C}$.

Next, we consider the case $\mathscr{C}=\mathscr{C}_w$.
For $X\in R(\beta)\text{-gmod}$, set
\[
W(X)\coloneq
\{\gamma\in Q_+\cap(\beta-Q_+)\mid e(\gamma,\beta-\gamma)X\neq0\}.
\]
By \cite[(2.5)]{strata}, we have $X\in\mathscr{C}_w$ if and only if
$W(X)\subseteq Q_+\cap wQ_-$.
Moreover, by \cite[Lemma~2.2]{strata},
\[
W(M\circ N)=W(M)+W(N).
\]
Since $0\in W(M),W(N)$, we obtain
\[
W(M),W(N)\subseteq W(M\circ N)\subseteq Q_+\cap wQ_-.
\]
Hence $M,N\in\mathscr{C}_w$.
\end{proof}
